
\documentclass{article}
\usepackage{colortbl}
\usepackage{makecell}
\usepackage{multirow}
\usepackage{supertabular}

\begin{document}

\newcounter{utterance}

\centering \large Interaction Transcript for game `codenames', experiment `abstract', episode 5 with reflexion{-}Qwen3.5{-}9B.
\vspace{24pt}

{ \footnotesize  \setcounter{utterance}{1}
\setlength{\tabcolsep}{0pt}
\begin{supertabular}{c@{$\;$}|p{.15\linewidth}@{}p{.15\linewidth}p{.15\linewidth}p{.15\linewidth}p{.15\linewidth}p{.15\linewidth}}
    \# & $\;$A & \multicolumn{4}{c}{Game Master} & $\;\:$B\\
    \hline

    \theutterance \stepcounter{utterance}  
    & & & \multicolumn{2}{p{0.3\linewidth}}{
        \cellcolor[rgb]{0.9,0.9,0.9}{
            \makecell[{{p{\linewidth}}}]{
                \texttt{\tiny{[GM$|$GM]}}
                {'hidden': {'team': ['priority', 'ethos'], 'innocent': ['urge', 'leniency', 'snobbery', 'hilarity'], 'opponent': [], 'assassin': []}, 'revealed': {'team': {'team': [], 'innocent': [], 'opponent': [], 'assassin': []}, 'opponent': {'team': [], 'innocent': [], 'opponent': [], 'assassin': []}}}
            }
        }
    }
    & & \\ \\

    \theutterance \stepcounter{utterance}  
    & & \multicolumn{4}{p{0.6\linewidth}}{
        \cellcolor[rgb]{0.9,0.9,0.9}{
            \makecell[{{p{\linewidth}}}]{
                \texttt{\tiny{[P1$\langle$GM]}}
                \texttt{You have attempted this same task before and did not succeed. These are the plans you wrote after those attempts.} \\
\texttt{Trial \#1: New plan: The previous attempt failed because the clue "attitude" was too broad and ambiguous, likely leading to incorrect guesses or allowing the opponent to guess a kill word. Specifically, "attitude" does not strongly connect to "desecration" or "mortality," weakening the semantic link for those targets. Next time, I will select a more precise, high{-}specificity clue that creates a tight semantic cluster for a subset of words (e.g., using "hubris" for "whim," "unbelief," and "sarcasm" to target arrogance{-}related concepts) rather than a generic term. I will also prioritize eliminating the "kill" word "manner" by ensuring my clue and targets do not loosely associate with it, and I will aim to guess the opponent's words first to disrupt their turn before revealing all my own targets.} \\
\\ 
\texttt{Trial \#2: New plan: The previous attempt failed because the clue "hubris" was not semantically related to "desecration" or "mortality," violating the core game rule that the clue must be a valid semantic link for the listed targets. Additionally, the format was correct, but the semantic validity was the primary error. Next time, I will select a clue that has a strong, direct definition{-}based connection to all listed targets (e.g., using "sin" for "desecration," "unbelief," "misbehavior," and "reluctance" if they fit, or grouping only tightly related words like "whim," "reluctance," and "misbehavior" under "impulse"). I will verify the semantic link for every single target word before submitting to ensure the clue is valid for the entire list.} \\
\\ 
\texttt{Trial \#3: New plan: The previous attempts failed because I included too many targets in a single clue, creating a semantic cluster that was either too broad (e.g., "death" for mortality, whim, sarcasm) or contained invalid links (e.g., "irony" for ethos/priority), causing the turn to end prematurely without clearing all words. Additionally, I failed to strategically target the opponent's remaining words to disrupt their turn. Next time, I will strictly limit each clue to a maximum of two tightly semantically linked words (e.g., CLUE: "stubbornness" TARGETS: "reluctance", "priority") to ensure 100\% validity. I will also prioritize guessing the opponent's specific remaining words ("exclusiveness", "inadvertence") in subsequent turns to force their turn end and gain extra opportunities to clear my own list.} \\
\\ 
\\ 
\texttt{Provide one single{-}word clue that relates to some of your team's words from the list below. You can choose to how many and to which words your clue relates to.} \\
\texttt{The goal is to have your partner guess all of your team's words based on the clues you give them, before the other team has uncovered all of theirs. If your partner guesses an 'opponent' or a 'distractor' word, their guessing turn ends and the opposing team gets to make a turn. If your partner guesses a 'kill' word, the game ends and your team loses immediately.} \\
\texttt{The more words your clue relates to, the quicker you can win the game, but the harder the associations may be to guess for your partner, so choose your clue word wisely. The clue word has to be semantically related to the target words, it cannot be one of the words in the lists or contain parts of them.} \\
\\ 
\texttt{Always give your single{-}word clue and your comma{-}separated list of related target words in the following format and make your answers as short as possible, never include any other text than is required in this form:} \\
\\ 
\texttt{CLUE: <WORD>} \\
\texttt{TARGETS: <TARGETS>} \\
\\ 
\texttt{Your team words are: priority, ethos.} \\
\texttt{Your opponent's team words are: .} \\
\texttt{Distractor words are: urge, leniency, snobbery, hilarity.} \\
\texttt{The kill words are: .} \\
            }
        }
    }
    & \\ \\

    \theutterance \stepcounter{utterance}  
    & \multicolumn{4}{p{0.6\linewidth}}{
        \cellcolor[rgb]{0.9,0.9,0.9}{
            \makecell[{{p{\linewidth}}}]{
                \texttt{\tiny{[P1$\rangle$GM]}}
                \texttt{CLUE: importance} \\
\texttt{TARGETS: priority, ethos} \\
            }
        }
    }
    & & \\ \\

    \theutterance \stepcounter{utterance}  
    & & & \multicolumn{2}{p{0.3\linewidth}}{
        \cellcolor[rgb]{0.9,0.9,0.9}{
            \makecell[{{p{\linewidth}}}]{
                \texttt{\tiny{[GM$|$GM]}}
                \texttt{importance} \\
            }
        }
    }
    & & \\ \\

    \theutterance \stepcounter{utterance}  
    & & & \multicolumn{2}{p{0.3\linewidth}}{
        \cellcolor[rgb]{0.9,0.9,0.9}{
            \makecell[{{p{\linewidth}}}]{
                \texttt{\tiny{[GM$|$GM]}}
                ['priority', 'ethos']
            }
        }
    }
    & & \\ \\

    \theutterance \stepcounter{utterance}  
    & & & \multicolumn{2}{p{0.3\linewidth}}{
        \cellcolor[rgb]{0.9,0.9,0.9}{
            \makecell[{{p{\linewidth}}}]{
                \texttt{\tiny{[GM$|$GM]}}
                {'word': 'priority', 'assignment': 'team'}
            }
        }
    }
    & & \\ \\

    \theutterance \stepcounter{utterance}  
    & & & \multicolumn{2}{p{0.3\linewidth}}{
        \cellcolor[rgb]{0.9,0.9,0.9}{
            \makecell[{{p{\linewidth}}}]{
                \texttt{\tiny{[GM$|$GM]}}
                {'word': 'ethos', 'assignment': 'team'}
            }
        }
    }
    & & \\ \\

    \theutterance \stepcounter{utterance}  
    & & & \multicolumn{2}{p{0.3\linewidth}}{
        \cellcolor[rgb]{0.9,0.9,0.9}{
            \makecell[{{p{\linewidth}}}]{
                \texttt{\tiny{[GM$|$GM]}}
                \texttt{opponent has won} \\
            }
        }
    }
    & & \\ \\

    \theutterance \stepcounter{utterance}  
    & \multicolumn{4}{p{0.6\linewidth}}{
        \cellcolor[rgb]{0.95,0.95,0.95}{
            \makecell[{{p{\linewidth}}}]{
                \texttt{\tiny{[P2$\rightarrow$P1]}}
                \texttt{Post{-}episode Reflexion (not sent to GM):} \\
\texttt{New plan: The previous attempt failed because the clue "importance" was too broad and semantically weak for "ethos," which specifically denotes a set of beliefs or character, rather than just general significance. This ambiguity likely allowed the opponent to guess a kill word or a distractor before "ethos" was secured. Next time, I will select a high{-}specificity clue that creates a tight, unambiguous semantic cluster for a small subset of words (e.g., using "core values" for "ethos" and "priority" if linked, or strictly "ethos" alone with a clue like "moral compass"). I will strictly limit each clue to one or two words to ensure 100\% validity and avoid triggering the opponent's turn prematurely.} \\
            }
        }
    }
    & & \\ \\

\end{supertabular}
}

\end{document}
